%% kapa: exam=1 round=34 year=2023 period=1 form=A subject=economics subject_ko=경제학원론 seq=5 number=45
%% source: Q-Net 감정평가사 기출문제 (cst00309) artlSeq=5211789 file=raw/1008/5211789/2023년도 제34회 감정평가사 1차시험 1교시.pdf
%% conversion: kordoc 4.24.51 -> tools/gen1.py (auto); review=proofread-sonnet; flags=none
\begin{question}{45}
소비자이론에 관한 설명으로 옳은 것을 모두 고른 것은?
\begin{framedtext}
ㄱ. 무차별곡선은 효용을 구체적인 수치로 표현할 수 있다는 가정 하에 같은 만족을 주는 점들을 연결한 것이다. \par
ㄴ. 상품의 특성에 따라 무차별곡선은 우상향 할 수도 있다. \par
ㄷ. 열등재이면서 대체효과보다 소득효과의 절대적 크기가 매우 클 경우 그 재화는 기펜재(Giffen goods)이다. \par
ㄹ. 유행효과(bandwagon effect)가 존재하면 독자적으로 결정한 개별수요의 수평적 합은 시장수요이다. \par
\end{framedtext}
\begin{choices}
  \item ㄱ, ㄴ
  \item ㄱ, ㄷ
  \item ㄱ, ㄹ
  \item ㄴ, ㄷ
  \item ㄴ, ㄹ
\end{choices}
\end{question}
\answer{4}
